% Part: sets-functions-relations
% Chapter: infinite
% Section: dedekinds-proof
%
\documentclass[../../../include/open-logic-section]{subfiles}

\begin{document}

\olfileid[pt]{sfr}{infinite}{dedekindsproof}

\olsection[``Prova'' de Dedekind]{``Prova'' de Dedekind da existência de um conjunto infinito}

Neste capítulo, apresentamos um tratamento em teoria dos conjuntos dos
números naturais, em termos de álgebras de Dedekind. Em
\olref[arith][ref]{sec}, refletimos sobre o significado filosófico
da aritmetização da análise (entre outras coisas).
Cumpre agora refletir sobre o significado do que aqui alcançamos.

Ao longo de \olref[sfr][arith][]{chap}, tomamos os números naturais como
dados e os usamos para construir os inteiros, os racionais e os reais,
de modo explícito. Neste capítulo, não demos uma construção explícita
dos números naturais. Mostramos apenas que,
\emph{dado qualquer conjunto infinito de Dedekind}, podemos definir um conjunto que se
comportará como queremos que~$\Nat$ se comporte.

Obviamente, pois, não podemos pretender ter respondido a uma questão
metafísica, como \emph{quais os objetos que são os números naturais}. Mas
isso é bom. Afinal, na \olref[sfr][arith][ref]{sec}, sublinhamos que
estaríamos errados se pensássemos na definição de
$\Real$ como conjunto de cortes de Dedekind como uma \emph{descoberta}, em vez de
uma estipulação conveniente. A observação crucial é que os cortes de Dedekind
exemplificam as propriedades matemáticas centrais dos números reais.
Assim também aqui: a observação crucial é que \emph{qualquer}
álgebra de Dedekind exemplifica as propriedades matemáticas centrais dos
números naturais. (De fato, Dedekind insistiu neste ponto ao provar
que todas as álgebras de Dedekind são \emph{isomorfas} (\citeyear[Theorems
132--3]{Dedekind1888}). Não surpreende, então, que muitos
``estruturalistas'' contemporâneos cite Dedekind como precursor.)

Além disso, mostramos como embutir a teoria dos números naturais
numa teoria ingênua simples dos conjuntos, que ainda permanece bastante
informal, mas que (aparentemente) não assume os números naturais como dados.
Assim, podemos estar no caminho de realizar o próprio projeto ambicioso de Dedekind,
que ele explicou assim:
\begin{quote}
	Na ciência nada que seja provável deve ser acreditado sem
	prova. Embora esta exigência pareça razoável, não posso considerá-la
	satisfeita sequer nos métodos mais recentes de lançar os fundamentos
	da ciência mais simples; a saber, aquela parte da lógica que trata
	da teoria dos números. Ao falar de aritmética (álgebra, análise) como
	mera parte da lógica, quero insinuar que considero o conceito de número
	totalmente independente das noções ou intuições de espaço e tempo---que
	o considero antes um produto imediato das leis puras do pensamento.
	\citep[preface]{Dedekind1888}
\end{quote}
A ideia audaciosa de Dedekind é esta. Acabamos de mostrar como construir os
números naturais usando só (ingenuamente) a teoria dos conjuntos. Em
\olref[sfr][arith][]{chap}, vimos como construir os reais dados
os números naturais e alguma teoria dos conjuntos. Logo, talvez, ``a aritmética
(álgebra, análise)'' acabe por ser ``meramente parte da lógica'' (no
sentido alargado que Dedekind dá à palavra ``lógica'').

Essa é a ideia. Mas esperemos um momento. A nossa construção de uma
álgebra de Dedekind (o nosso substituto dos números naturais) é
condicional à existência de um conjunto infinito de Dedekind. (Basta olhar
para trás, para \olref[sfr][infinite][dedekind]{thm:DedekindInfiniteAlgebra}.)
A menos que a existência de um conjunto infinito de Dedekind possa ser estabelecida via
``lógica'' ou ``as leis puras do pensamento'', o projeto estagna.

Então, \emph{pode} a existência de um conjunto infinito de Dedekind ser estabelecida
pelas ``leis puras do pensamento''? Eis o esforço de Dedekind:
\begin{quote}
  O meu próprio reino de pensamentos, O.s.\@, a totalidade $S$ de todas as coisas que
  podem ser objetos do meu pensamento, é infinita. Pois, se $s$ designa um
  elemento de~$S$, então o pensamento $s'$ de que~$s$ pode ser objeto do
  meu pensamento é ele próprio um elemento de~$S$. Se considerarmos isto como uma
  imagem $\phi(s)$ do elemento~$s$, então \dots~$S$ é infinito de Dedekind,
  como queríamos ter provado.
	\citep[\S66]{Dedekind1888}
\end{quote}
Isto é algo de espantoso de encontrar no meio de um livro
que em grande parte consiste em provas matemáticas altamente rigorosas. Duas
observações valem a pena.

Primeiro: esta ``prova'' mal tem o que hoje reconheceríamos como caráter
``matemático''. Fala de objetos psicológicos
(pensamentos), e apenas de objetos \emph{possíveis}.

Segundo: pelo menos tal como apresentamos as álgebras de Dedekind, esta
``prova'' tem uma lacuna técnica direta. Se o argumento de Dedekind é bem-sucedido,
estabelece só que há infinitas
coisas (nomeadamente, infinitos pensamentos). Mas Dedekind
também precisa de nos dar uma razão para ver~$S$ como um único \emph{conjunto},
com infinitos !!{element}s, em vez de pensar em~$S$ como
\emph{algumas coisas} (no plural).

O fato de Dedekind não ver aqui uma lacuna pode sugerir que o seu
uso da palavra ``totalidade'' não corresponde exatamente ao \emph{nosso} uso
da palavra ``conjunto''.\footnote{De fato, temos outras razões para pensar
que não correspondia; ver \citet[p.~23]{Potter2004}.} Mas isto não seria
demasiado surpreendente. O projeto que perseguimos nos últimos dois
capítulos---uma ``construção'' dos naturais e, a partir deles, uma
``construção'' dos inteiros, reais e racionais---foi todo ele
feito de modo ingênuo. Nos apoiamos neste conjunto, ou naquele
conjunto, sempre que precisamos, sem enunciar muitos princípios gerais
sobre exatamente que conjuntos existem, e quando. Mas sabemos
que precisamos de \emph{alguns} princípios gerais, pois caso contrário caímos
no Paradoxo de Russell.

Chegou a hora de superarmos a nossa ingenuidade.

\end{document}
